\section[Rédaction Bourbaki \redac{596} : cinq énoncés sur les modules]{Rédaction Bourbaki \redac{596} : cinq énoncés sur les modules à corriger\protect\verifie{Redaction596Modules}}\label{sec:bourbaki-596-modules}

Les premiers paragraphes de la rédaction \redac{596} (section~\ref{sec:bourbaki-596}), sur les modules de longueur finie, la théorie de Morita et les anneaux simples, contiennent cinq énoncés faux ou imprécis tels qu'ils sont tapés, en plus de ceux de la section~\ref{sec:bourbaki-596-hypotheses}. La transcription reste locale ; on n'en cite aucun passage. On note $A_s$ (resp.\ $A_d$) l'anneau $A$ vu comme module à gauche (resp.\ à droite) sur lui-même et $A^{\circ}$ l'anneau opposé.

\emph{§~4, n\textsuperscript{o}\,3, exemple~1 (p.~54 du fac-similé).} Soient $S = A_s^n$ et $D = M_n(A^{\circ}) = \operatorname{End}_A(S)$. Le $D$-module $S$ serait artinien (resp.\ noethérien) si et seulement si le $A$-module $A_s$ l'est. Or les sous-$D$-modules de $S$ correspondent aux idéaux à \emph{droite} de $A$ : la condition porte sur $A_d$. Les autres assertions de l'exemple (semi-simple, indécomposable, simple) sont justes, parce que ces propriétés de $A_s$ et de $A_d$ coïncident.

\emph{§~4, remarque~4 (en astérisques, p.~55).} Les classes d'autoéquivalences de la catégorie des $A$-modules forment un groupe ; pour $A$ commutatif non nul, il serait isomorphe au groupe $\operatorname{Pic}(A)$ des classes de modules projectifs de rang~$1$. C'est vrai des autoéquivalences $A$-linéaires. En général un automorphisme $\sigma$ de l'anneau $A$ donne une autoéquivalence, la restriction des scalaires par $\sigma$, et le groupe est $\operatorname{Pic}(A) \rtimes \operatorname{Aut}(A)$.

\emph{§~5, n\textsuperscript{o}\,6, définition~5 et proposition~16\,a) (pp.~72--73).} Soient $B$ un sous-anneau simple de $A$ et $M \neq 0$ un $A$-module de longueur finie sur $B$. Le commutant $A'$ du $A$-module $M$ serait un anneau simple. La démonstration applique le corollaire~1 du théorème~1 du n\textsuperscript{o}\,4, qui porte sur des modules semi-simples isotypiques ; il faut donc $A$ simple, ce que le numéro ne suppose pas.

\emph{§~2, n\textsuperscript{o}\,2, exemple~5 (p.~33).} Sur un anneau principal qui n'est pas un corps, un module de torsion dont les diviseurs élémentaires sont « de multiplicité~$1$ » serait semi-simple. Si le mot veut dire que chaque diviseur élémentaire apparaît une seule fois, l'énoncé est faux, et la condition n'est pas non plus nécessaire. La condition voulue est que chaque diviseur élémentaire soit un élément premier, c'est-à-dire d'exposant~$1$.

\emph{§~1, n\textsuperscript{o}\,1, exemple~3 (p.~9).} Le $\mathbf Z$-module $\mathbf Z/n\mathbf Z$ serait indécomposable si et seulement si $n = p^m$, avec $p$ premier et $m \geq 1$. C'est vrai pour $n \geq 1$, faux pour $n = 0$.

\begin{encadre}
\begin{contrexemple}[à cinq énoncés des §§~1, 2, 4 et~5 de la rédaction \redac{596}]
\begin{enumerate}
\item[a)] Soit $T = \left(\begin{smallmatrix} \mathbf R & \mathbf R \\ 0 & \mathbf Q \end{smallmatrix}\right)$, sous-anneau de $M_2(\mathbf R)$, et $n = 1$, de sorte que $D = T^{\circ}$. Le $T$-module $T_s$ est artinien et noethérien ; le $T^{\circ}$-module $T_s$, dont les sous-modules sont les idéaux à droite de $T$, n'est ni artinien ni noethérien.
\item[b)] Soit $\sigma$ la conjugaison de $\mathbf C$. La restriction des scalaires par $\sigma$ est une autoéquivalence de la catégorie des $\mathbf C$-espaces vectoriels qui n'est pas isomorphe à l'identité, alors que $\operatorname{Pic}(\mathbf C)$ est trivial.
\item[c)] Soient $A = \mathbf Q[\varepsilon]$, $\varepsilon^2 = 0$, $B = \mathbf Q$ et $M = A$, de longueur~$2$ sur $B$. Le commutant $\operatorname{End}_A(A) \cong A$ n'est pas un anneau simple.
\item[d)] Le $\mathbf Z$-module $\mathbf Z/4$, de seul diviseur élémentaire $4$, n'est pas semi-simple ; le $\mathbf Z$-module $\mathbf Z/2 \times \mathbf Z/2$, où le diviseur $2$ apparaît deux fois, l'est.
\item[e)] Le $\mathbf Z$-module $\mathbf Z/0\mathbf Z = \mathbf Z$ est indécomposable, et $0$ n'est pas de la forme $p^m$.
\end{enumerate}
\end{contrexemple}
\end{encadre}

\begin{proof}
a) Le sous-anneau diagonal $C = \mathbf R \times \mathbf Q$ agit sur $T$ par multiplication à gauche, et $T$ est engendré sur $C$ par $e_{11}$, $e_{12}$, $e_{22}$ : on a $\left(\begin{smallmatrix} x & y \\ 0 & q \end{smallmatrix}\right) = (x, 0)\,e_{11} + (y, 0)\,e_{12} + (0, q)\,e_{22}$. Comme $C$ est artinien et noethérien, le $C$-module $T$ l'est ; ses sous-$T$-modules à gauche sont des sous-$C$-modules, donc $T_s$ est artinien et noethérien. À droite, pour tout sous-$\mathbf Q$-espace vectoriel $V$ de $\mathbf R$, l'ensemble $\left(\begin{smallmatrix} 0 & V \\ 0 & 0 \end{smallmatrix}\right)$ est un idéal à droite, puisque $\left(\begin{smallmatrix} 0 & v \\ 0 & 0 \end{smallmatrix}\right)\left(\begin{smallmatrix} x & y \\ 0 & q \end{smallmatrix}\right) = \left(\begin{smallmatrix} 0 & qv \\ 0 & 0 \end{smallmatrix}\right)$. L'application $V \mapsto \left(\begin{smallmatrix} 0 & V \\ 0 & 0 \end{smallmatrix}\right)$ est un plongement d'ensembles ordonnés ; si $T$ était artinien ou noethérien à droite, $\mathbf R$ serait un $\mathbf Q$-espace vectoriel artinien ou noethérien, donc de dimension finie, donc dénombrable. Le transposé $\left(\begin{smallmatrix} \mathbf Q & \mathbf R \\ 0 & \mathbf R \end{smallmatrix}\right)$, isomorphe à $T^{\circ}$, sépare les deux conditions dans l'autre sens.

b) Soit $\alpha$ un isomorphisme de la restriction des scalaires par $\sigma$ vers l'identité, et $f : \mathbf C \to \mathbf C$ sa composante en $\mathbf C$. La linéarité donne $f(\sigma(c)\,x) = c\,f(x)$ ; la naturalité, appliquée aux multiplications par $b$, qui sont $\mathbf C$-linéaires, donne $f(bx) = b\,f(x)$. Donc $\sigma(c)\,f(1) = f(\sigma(c)) = c\,f(1)$, et $f(1)$ est inversible puisque $f$ est surjective ; d'où $\sigma(c) = c$ pour tout $c$, ce qui est faux pour $c = i$. L'argument vaut pour tout anneau commutatif et tout automorphisme $\sigma \neq \mathrm{id}$. Enfin $\mathbf C$ est local, donc tout module projectif de rang~$1$ est libre.

c) L'anneau $\operatorname{End}_A(A)$ est isomorphe à $A^{\circ} = A$. Un anneau commutatif simple est un corps ; or $\varepsilon \neq 0$ et $\varepsilon^2 = 0$.

d) Soit $N$ un supplémentaire de $P = 2\mathbf Z/4\mathbf Z$ dans $\mathbf Z/4$. Écrivons $1 = p + x$, avec $p \in P$ et $x \in N$. Comme $2p = 0$, on a $2x = 2 \in P \cap N$, et $2 \neq 0$ : contradiction. Le groupe $\mathbf Z/2 \times \mathbf Z/2$ est un espace vectoriel sur $\mathbf Z/2$, et ses sous-$\mathbf Z$-modules sont ses sous-espaces.

e) Si $\mathbf Z = P \oplus Q$ avec $a \in P$ et $b \in Q$ non nuls, alors $ab$ appartient aux idéaux $P$ et $Q$, et $ab \neq 0$.
\end{proof}

\begin{sloppypar}
\noindent\emph{Vérification en Lean.} a) est \texttt{isArtinianRing\_T}, \texttt{isNoetherianRing\_T}, \texttt{not\_isArtinian\_op} et \texttt{not\_isNoetherian\_op}, réunis dans \texttt{example1\_fails} ; le plongement des sous-espaces est \texttt{cornerEmb} ; b) est \texttt{eq\_of\_iso\_id}, pour tout anneau commutatif, puis \texttt{restrictScalars\_conj\_not\_iso\_id} et \texttt{pic\_complex\_subsingleton} ; c) est \texttt{finite\_dualNumber} et \texttt{prop16a\_fails} ; d) est \texttt{not\_isSemisimpleModule\_zmod\_four} et \texttt{isSemisimpleModule\_zmod\_two\_sq} ; e) est \texttt{int\_indecomposable} et \texttt{zero\_ne\_prime\_pow}, dans \href{\leanblobmain/Redaction596Modules.lean}{\texttt{Redaction596\allowbreak{}Modules.\allowbreak{}lean}}. Seul le cas $n = 1$ de a) est formalisé ; la structure du groupe des autoéquivalences dans b) ne l'est pas.
\end{sloppypar}

\litterature{les énoncés justes sont classiques : Bass, \emph{Algebraic K-Theory} (1968), ch.~II, pour Morita et le groupe de Picard ; Lam, \emph{A First Course in Noncommutative Rings} (2\textsuperscript{e} éd., 2001), §~1, pour les anneaux triangulaires artiniens d'un seul côté. Le texte imprimé de 2012 n'a pas été consulté : on ne sait pas si ces énoncés y sont corrigés.}
